\section{The $q$-infinitesimal C1.1 source at both small testing angles}
\label{sec:r7-c1-q-inf-local}
\begin{theorem}
On the original C1.1 source with $q=k^2\eta$ infinitesimal, the entire
original parent is tensor SOS whenever both independent testing
half-angles $u,v$ are infinitesimal. All original opening rates and
$\beta=0$ are retained. The boundary $D=0$ is empty here because
$D_b$ has residue $1/2$.
\end{theorem}
This statement includes every state; it does not yet assert the
ordinary or unbounded $v$ testing charts.

Use the preceding source notation and put $c=\sqrt q$, $e=c/g$,
$h=H/R$, and $E_*=h-v$. Work first with $Z=1$ and $u,v,h>0$.
Define
\[
 A=ch,\quad B=cL/R,\quad C=uv/b,\quad D=qg,\quad
 \Gamma=(ABCD)^{1/4},\quad d_i=A_i-\Gamma.
\]
Then
\[
 \Gamma^2=qhv,\qquad BC=e v^2h,\qquad
 u h=uE_*+bC.
 \tag{QI1}
\]
Here $h=u(W+uY)/(Rb^2)$ is separated into a positive pure-source
factor and a positive pure-testing factor. Every radical below is
therefore legal; no mixed inverse or order comparison is used.
The exact original inequality $\beta\ge0$ gives $t=O(e^2)$, while
$t$ is bounded. Hence $t/e$ and $\sqrt{t/e}$ are bounded.

\paragraph{Exact scalar, with the $D$ packet restored.}
Write $U=u^2,V=v^2$, $D_u=1+U,D_v=1+V$, and
$a_t=1-\rho_0+(t/R)V$. At both state normals zero the entire original
scalar is
\[
 a h^2+b_*hv+c_*v^2+P_L L^2+2D_uC^2+4\alpha_b g^2D_u,
 \tag{QI2}
\]
where
\begin{align*}
 a={}&4D_vR^2\{c_b(1-\rho_0)^2+a_b\rho_0^2/D_u\}
 +2b^2R^2\{a_t^2+V\rho_0^2/D_u\}\\
 &+4\alpha_b g^2b^4VR^2\rho_0^2/D_u
 +4D_bb^4R^2a_t^2,\\
 b_*={}&-4R\{D_u(1+tV/R)-U\rho_0\}
 -8\alpha_b g^2b^2R\rho_0-8D_bb^2Ra_tD_u,\\
 c_*={}&4D_bD_u^2.
\end{align*}
All these rational functions have bounded source coefficient arrays;
only $D_u$ and positive pure-source units occur in their denominators.
The $r$ openings cancel exactly as in (SC12).
In particular
\[
 D_b=\frac{N_D}{2JP_d},\qquad
 N_D=\frac{K(1+qg^2)(1+b^2r)-JP_d}{b^2}
\]
has a polynomial numerator: the numerator before division vanishes
identically at $b^2=0$, where $P_d=1+qg^2$. This is a finite monic
division, so it introduces no unbounded source coefficient array.
Also $\alpha_b/q^2$, $h_2/q$, and $h_b/q$ have explicit bounded
arrays from their source formulas. These observations justify the
$q$ coquotients below uniformly through both source openings.

Use formal specialization at $q=0$ only to construct finite polynomial
coquotients; it is not a new physical source. Put $z=b^2$,
$J_0=1-z(1-r)$, and $S=1-r+tV$. The three specialized coefficients are
\begin{align*}
 a_0={}&2D_v\{(1+zr)(1-r)+r/D_u\}
       +2z\{S^2+Vr^2/D_u\}+2z^2S^2/J_0,\\
 b_0={}&-4\{D_u(1+tV)-Ur\}-4zSD_u/J_0,\\
 c_0={}&2D_u^2/J_0.
\end{align*}
Define $\widetilde b_0=b_0+4tD_uV$, and the bounded finite
coquotients
\[
 X=a_0+\widetilde b_0+c_0=U X_U+V X_V,
 \qquad Y=-\widetilde b_0-2c_0=U Y_U+V Y_V.
 \tag{QI3}
\]
For definiteness, first subtract the value at $U=0$, divide by $U$,
and then divide that value by $V$. Their two $V$ coquotients are
\[
 X_V=2+2zr(1-2t)+\frac{2z(1+zr)t^2V}{J_0},
 \qquad Y_V=\frac{4zt}{J_0}.
 \tag{QI4}
\]
Thus $X_V=2+\varepsilon_V$, $\varepsilon_V\in B_{\rm Inf}$;
$X_U,Y_U$ have bounded arrays. Likewise write
$a=a_0+qa_1$, $b_*=b_0+qb_1$, $c_*=c_0+qc_1$ by finite numerator
division. All $a_1,b_1,c_1$ have bounded arrays. Direct source
specialization gives
\[
 \overline{a_1-c_1}=2,\qquad
 \overline{b_1+2c_1}=-8,\qquad
 \overline{R^2P_L/q}=2,\qquad
 \overline{4\alpha_bD_u/q^2}=2.
 \tag{QI5}
\]

The following is the exact scalar identity to be allocated:
\begin{align*}
 (\mathrm{QI2})={}&(c_0+qc_1)E_*^2+2D_uC^2
 +X_U(uE_*+bC)^2+Y_Uu(uE_*+bC)E_*\\
 &+X_V(vh)^2+Y_Vv(vh)E_*
 -4(t v/e)D_uBC\\
 &+(a_1-c_1)A^2+(b_1+2c_1)\Gamma^2
 +(R^2P_L/q)B^2+(4\alpha_bD_u/q^2)D^2.
 \tag{QI6}
\end{align*}
It uses only (QI1) and $qv^2=2\Gamma^2-A^2+qE_*^2$.
In particular the apparently dangerous negative $t$ term is already
an infinitesimal-coefficient $BC$ term, since $v$ is infinitesimal.

\paragraph{Two extra contact squares and their finite first jets.}
Set
\[
 F_*=vh-gD,\quad
 R_F=\frac{\sqrt{g\Gamma}}{c}(\Gamma-D),\quad
 V_0=cv,\quad R_E=\sqrt\Gamma(\sqrt A-\sqrt{V_0}).
\]
All four roots vanish at actual contact. Exact multiplication gives
\begin{align*}
 (vh)^2&=F_*^2+2R_F^2+2gF_*d_4+g^2D^2-4g^2\Gamma d_4,\\
 c\Gamma E_*&=2\Gamma d_1-R_E^2,\\
 \Gamma F_*&=-2g\Gamma d_4+R_F^2/g.
 \tag{QI7}
\end{align*}
The root $R_E$ receives a strict real square pivot from $E_*^2$:
\[
 qE_*^2=(A-V_0)^2
       =4R_E^2+(\sqrt A-\sqrt{V_0})^4.
 \tag{QI8}
\]
Reserve this fraction $q$ of the strict $E_*$ pivot; the remaining
coefficient still has positive real residue. The leading $2(vh)^2$
in (QI6), by (QI7), gives strict $F_*$ and $R_F$ pivots, with only
infinitesimal quadratic corrections and the printed first-jet term.
The remaining leading angular form is $2(\sum A_i^2-4\Gamma^2)$,
with the strict finite toric square list from (SC16).

For precision, the only additional bare-$\Gamma$ normal couplings
left by (QI6), apart from ordinary toric first jets, have the form
$\Gamma E_*a_*+\Gamma F_*b_*'$ with
\[
 a_*=bu(2X_U+Y_U)+Y_Vvg,
 \qquad b_*'=2\varepsilon_Vg.
\]
At contact $u=bc,v=cg$ these satisfy
\[
 a_{*0}/c=b^2(2X_{U0}+Y_{U0})+Y_{V0}g^2\in B_{\rm Inf},
 \qquad b_{*0}'/g=2\varepsilon_{V0}\in B_{\rm Inf}.
\]
Their differences from contact have finite bounded monic Taylor
arrays. The conductors (SC18) allocate those differences against
$E_*,F_*$ with infinitesimal coefficients. Apply (QI7) to their
constant terms. What remains is a bounded infinitesimal toric
quadratic plus $\Gamma\sum\mu_i d_i+\Gamma^2\sigma$.
The exact scalar value and three original first jets imply its
contact value and first-jet balances. The same two finite divisions
as in (SC17)--(SC20), or the Hadamard contact-lift identity, allocate
it on the strict toric root list. No infinite series or unsupported
ideal-membership assertion is used: (QI3), the $q$ coquotients,
and the two displayed contact coefficient ratios are the finite
certificates needed for the additional roots.
Every coefficient function in these Taylor steps depends only on
$u,v$. Thus only the first two conductors in (SC18)--(SC20) are
used; the third state-coordinate jet supplies the inactive Hadamard
balance. No bound on $(W+uY)/R$ after $Z$ normalization is assumed.

\paragraph{Both complete state normals.}
Use $\widehat N=\sqrt{H_0}N$, $H_0=4h_2/R+4h_b/t$, as in the
preceding uniform-normal lemma. Add back the centered $D$ packet
from (SC10), including its normals. Its scalar root is exactly
\[
 J_*=(T+b^2tV)h-D_uv
 =E_*+b^2c(1-tg^2)A-buC+b^2tvF_*+b^2tvgD.
 \tag{QI9}
\]
Its full root is $J_*+ubvgN-b^2N_y/\sqrt s$.
All additional hard coefficients are infinitesimal after the
$\widehat N$ normalization.

For the first $N$ normal in (SC14), its $H$ term is a bounded
infinitesimal multiple of $u h=uE_*+bC$; its other terms are
infinitesimal multiples of $B,C$. For the harmonic normal, multiply
its $H$ coefficient by $\sqrt s$. After the exact opening
cancellations, its value at $q=U=V=0$ is zero. Finite numerator
coquotients therefore give
\[
 \sqrt s\,[H\text{-coefficient in }\mathcal L_y]
       =qK_q+UK_U+VK_V
 \tag{QI10}
\]
with bounded arrays. Use
\[
 qH=RcA,\quad u^2H=Ru(uE_*+bC),\quad
 v^2H=Rv(F_*+gD).
\]
The two remaining terms of $\mathcal L_y$ have infinitesimal
coefficients on $C,D$ since $\alpha_b/q^2$ is bounded. Thus both
complete first normals, including (QI9), are infinitesimal
combinations of $E_*,F_*,A,B,C,D$. Their exact contact zeros make
the coefficient sum of $A,B,C,D$ zero at contact. One finite monic
division and (SC18) allocate these on the same strict square list.
The hard block has residue $\operatorname{diag}(1,4)$ and all its
other entries have zero residue. This proves a bounded positive
residue Gram for the entire original parent, after restoring all
retained positive packets.

\paragraph{States and signs.}
The restriction $Z=0$ is positive by the original pure-testing state
restriction. For $Z\ne0$, divide the state by the pure-testing $Z$;
its size is immaterial, since every coefficient array above depends
only on $u,v$. On the negative-core chart, exact sign gauges give
$u,v>0$ and the inherited pure-testing Schur identity forces
$W+uY>0$. Hence $h>0$, as required. The other core-sign chart and
literal testing axes are direct original exits. Complex states
split into real and imaginary parts. This completes the asserted
small-$u$, small-$v$ theorem.
